\chapter{Simple Linear Regression}
\label{ch:simple-linear-regression}

Here are six file uploads. For each one we record the file size $x$ (in MB) and
the upload time $y$ (in seconds).

\begin{center}
\begin{tabular}{lcccccc}
\toprule
upload $i$ & 1 & 2 & 3 & 4 & 5 & 6 \\
\midrule
size $x_i$ (MB) & 1 & 2 & 3 & 4 & 5 & 6 \\
time $y_i$ (s) & 3 & 5 & 4 & 7 & 8 & 9 \\
\bottomrule
\end{tabular}
\end{center}

Bigger files take longer, but not perfectly: the 3 MB file was faster than the
2 MB file. Chapters~9 to~11 compared groups: version A against version B, one
mean against another. Here $x$ is not a group label but a number that can take
any value, and we want to know how $y$ changes as $x$ slides. Chapter~4
measured the strength of such a relationship with the correlation $r$, a
number between $-1$ and $1$ that says how tightly the points hug a line. This
chapter goes further. We fit a line, use it to predict,
say how much of the variation it explains, and put error bars on its slope.
That is \textbf{simple linear regression}: one predictor $x$, one response $y$.

\section{The model}
\label{sec:slr-model}

\begin{definition}[Simple linear regression model]
The response is a straight-line function of $x$ plus noise:
\[
  y_i = \beta_0 + \beta_1 x_i + \eps_i, \qquad i = 1,\dots,n,
\]
where the errors $\eps_i$ are independent, have mean $0$, and have the same
variance $\sigma^2$. Often we also assume $\eps_i \sim \Normal(0,\sigma^2)$.
The unknown parameters are the \textbf{intercept} $\beta_0$, the \textbf{slope}
$\beta_1$, and the \textbf{noise level} $\sigma$.
\end{definition}

The model says: at each fixed $x$, the response has a distribution. Its mean
sits on the line $\beta_0+\beta_1 x$, and its spread is $\sigma$ no matter
where you are. Figure~\ref{fig:slr-model} draws this picture. The line is the
\emph{average} upload time for a file of a given size. Individual uploads
scatter around it because of network jitter, server load, and everything else
we did not measure.

\begin{figure}[H]
\centering
\begin{tikzpicture}
\begin{axis}[
  width=11cm, height=6.5cm, axis lines=left,
  xmin=0, xmax=7.5, ymin=0, ymax=13,
  xlabel={file size $x$ (MB)}, ylabel={upload time $y$ (s)},
  clip=false,
]
\addplot[sourceblue, very thick, domain=0:7.2] {1.8+1.2*x};
\node[sourceblue, anchor=west] at (axis cs:7.2,10.44) {$\E[y\given x]=\beta_0+\beta_1x$};
\foreach \xo in {1.5,3.5,5.5} {
  \addplot[densityred, thick, domain=-2.4:2.4, samples=60]
    ({\xo+1.1*exp(-x^2/1.28)}, {1.8+1.2*\xo+x});
  \addplot[densityred, dashed] coordinates {(\xo,{1.8+1.2*\xo-2.4}) (\xo,{1.8+1.2*\xo+2.4})};
}
\addplot[only marks, mark=*, mark size=1.8pt, black] coordinates {
 (1.5,2.9) (1.5,4.4) (1.5,3.6) (3.5,6.9) (3.5,5.2) (3.5,6.3) (5.5,8.0) (5.5,9.6) (5.5,8.8)};
\node[densityred, anchor=west, font=\small] at (axis cs:6.7,6.3) {same spread $\sigma$};
\draw[densityred, -{Stealth}] (axis cs:6.7,6.3) -- (axis cs:6.4,7.6);
\end{axis}
\end{tikzpicture}
\caption{The regression model. At every $x$ the response is a random variable
centred on the line, with the same spread $\sigma$ (red bells). Observed points
(black) are single draws from these bells.}
\label{fig:slr-model}
\end{figure}

\begin{intuitionbox}
Read $\beta_1$ as ``how much the \emph{average} $y$ moves when $x$ goes up by
one unit.'' Read $\sigma$ as ``how far a typical point lands from that average.''
Regression estimates both. A small slope with small $\sigma$ can be a clearer
signal than a large slope buried in large $\sigma$.
\end{intuitionbox}

We never see $\beta_0,\beta_1,\sigma$. We see $n$ points. Our job is to compute
estimates $b_0, b_1, s$ from the data and then say how far off they might be.
(Many texts, and Chapters~13--15, write $\hat\beta_0, \hat\beta_1$ for $b_0, b_1$;
the hat means ``estimated from data''.) The fitted line is written
\[
  \hat y = b_0 + b_1 x .
\]
The hat means ``predicted value.'' For the $i$-th point, $\hat y_i = b_0 + b_1 x_i$.

\section{Residuals and the least squares criterion}
\label{sec:slr-ls}

\begin{definition}[Residual]
The \textbf{residual} of point $i$ is the observed value minus the predicted value:
\[
  e_i = y_i - \hat y_i = y_i - (b_0 + b_1 x_i).
\]
Geometrically it is the signed \emph{vertical} distance from the point to the line.
\end{definition}

A positive residual means the point lies above the line (the model
under-predicted). Negative means below. Figure~\ref{fig:slr-residuals} shows
the residuals of our six uploads for the line we are about to derive.

\begin{figure}[H]
\centering
\begin{tikzpicture}
\begin{axis}[
  width=11cm, height=6.5cm, axis lines=left,
  xmin=0, xmax=7, ymin=1, ymax=10.5,
  xlabel={file size $x$ (MB)}, ylabel={upload time $y$ (s)},
]
\addplot[sourceblue, very thick, domain=0.3:6.7] {1.8+1.2*x};
\addplot[residgray, thick] coordinates {(2,5) (2,4.2)};
\addplot[residgray, thick] coordinates {(3,4) (3,5.4)};
\addplot[residgray, thick] coordinates {(4,7) (4,6.6)};
\addplot[residgray, thick] coordinates {(5,8) (5,7.8)};
\addplot[only marks, mark=*, mark size=2.2pt, black] coordinates {(1,3) (2,5) (3,4) (4,7) (5,8) (6,9)};
\node[anchor=east, font=\small] at (axis cs:1.95,4.6) {$e_2=+0.8$};
\node[anchor=west, font=\small] at (axis cs:3.05,4.5) {$e_3=-1.4$};
\node[anchor=east, font=\small] at (axis cs:3.95,7.3) {$e_4=+0.4$};
\addplot[only marks, mark=square*, mark size=2.6pt, quotientorange] coordinates {(3.5,6)};
\node[quotientorange, anchor=north west, font=\small] at (axis cs:3.55,5.95) {$(\xbar,\ybar)=(3.5,6)$};
\node[sourceblue, anchor=south east, font=\small] at (axis cs:6.6,9.7) {$\hat y = 1.8+1.2x$};
\end{axis}
\end{tikzpicture}
\caption{Residuals are vertical segments from each point to the fitted line.
Points 1 and 6 sit exactly on the line ($e=0$). The orange square is the
centroid $(\xbar,\ybar)$; the least squares line always passes through it.}
\label{fig:slr-residuals}
\end{figure}

\begin{trapbox}
Residuals are \emph{vertical}, not perpendicular. We are predicting $y$ from
$x$, so the error that matters is the error in $y$ at the given $x$. Minimising
perpendicular distances gives a different line (total least squares, or the
first principal component), which answers a different question.
\end{trapbox}

Which line is ``best''? We want residuals small as a whole. Their plain sum is
useless: big positive and big negative residuals cancel. The standard choice is
to square them and add.

\begin{definition}[Least squares]
The \textbf{least squares} line is the choice of $(b_0,b_1)$ that minimises the
sum of squared errors
\[
  \mathrm{SSE}(b_0,b_1) = \sum_{i=1}^n \bigl(y_i - b_0 - b_1 x_i\bigr)^2 .
\]
\end{definition}

Squaring has a literal picture. Each squared residual is the area of a square
whose side is the residual. Least squares finds the line with the smallest total
square area (Figure~\ref{fig:slr-squares}).

\begin{figure}[H]
\centering
\begin{tikzpicture}
\begin{axis}[
  name=left, width=6.4cm, height=6.4cm, axis lines=left,
  axis equal image, xmin=0, xmax=9.5, ymin=1, ymax=10.5,
  xlabel={$x$}, ylabel={$y$}, title style={font=\small}, title={\small flat line $\hat y=6$: $\mathrm{SSE}=28$},
  xtick={0,2,4,6,8}, ytick={2,4,6,8,10},
]
\addplot[trapred, very thick, domain=0.3:9] {6};
\draw[trapred, fill=trapred, fill opacity=0.18] (axis cs:1,3) rectangle (axis cs:4,6);
\draw[trapred, fill=trapred, fill opacity=0.18] (axis cs:2,5) rectangle (axis cs:3,6);
\draw[trapred, fill=trapred, fill opacity=0.18] (axis cs:3,4) rectangle (axis cs:5,6);
\draw[trapred, fill=trapred, fill opacity=0.18] (axis cs:4,7) rectangle (axis cs:5,6);
\draw[trapred, fill=trapred, fill opacity=0.18] (axis cs:5,8) rectangle (axis cs:7,6);
\draw[trapred, fill=trapred, fill opacity=0.18] (axis cs:6,9) rectangle (axis cs:9,6);
\addplot[only marks, mark=*, mark size=1.8pt, black] coordinates {(1,3) (2,5) (3,4) (4,7) (5,8) (6,9)};
\end{axis}
\begin{axis}[
  at={(left.east)}, anchor=west, xshift=0.6cm,
  width=6.4cm, height=6.4cm, axis lines=left,
  axis equal image, xmin=0, xmax=9.5, ymin=1, ymax=10.5,
  xlabel={$x$}, title style={font=\small}, title={\small least squares: $\mathrm{SSE}=2.8$},
  xtick={0,2,4,6,8}, ytick={2,4,6,8,10},
]
\addplot[sourceblue, very thick, domain=0.3:7] {1.8+1.2*x};
\draw[sourceblue, fill=sourceblue, fill opacity=0.25] (axis cs:2,5) rectangle (axis cs:2.8,4.2);
\draw[sourceblue, fill=sourceblue, fill opacity=0.25] (axis cs:3,4) rectangle (axis cs:4.4,5.4);
\draw[sourceblue, fill=sourceblue, fill opacity=0.25] (axis cs:4,7) rectangle (axis cs:4.4,6.6);
\draw[sourceblue, fill=sourceblue, fill opacity=0.25] (axis cs:5,8) rectangle (axis cs:5.2,7.8);
\addplot[only marks, mark=*, mark size=1.8pt, black] coordinates {(1,3) (2,5) (3,4) (4,7) (5,8) (6,9)};
\end{axis}
\end{tikzpicture}
\caption{Squared residuals drawn as actual squares. The flat line at the mean
has total square area $28$; the least squares line cuts it to $2.8$. No other
line does better.}
\label{fig:slr-squares}
\end{figure}

\begin{whybox}
Why squares and not absolute values? Three reasons. (1) Squares make
$\mathrm{SSE}$ a smooth bowl in $(b_0,b_1)$, so calculus gives a closed-form
answer. (2) The answer is built from means, variances and covariances, the same
quantities as Chapters~2 and~4. (3) If the noise is Gaussian, least squares is
exactly maximum likelihood (Section~\ref{sec:slr-mle}). Minimising absolute
residuals instead (least absolute deviations) is more robust to outliers, just
as the median is more robust than the mean, but has no closed form.
\end{whybox}

\section{Deriving the least squares line}
\label{sec:slr-derive}

We need some shorthand. For data $(x_1,y_1),\dots,(x_n,y_n)$ define
\begin{align*}
  S_{xx} &= \sum_{i=1}^n (x_i-\xbar)^2, &
  S_{yy} &= \sum_{i=1}^n (y_i-\ybar)^2, &
  S_{xy} &= \sum_{i=1}^n (x_i-\xbar)(y_i-\ybar).
\end{align*}
These are the un-normalised versions of the variances and covariance:
$s_x^2 = S_{xx}/(n-1)$, $s_y^2 = S_{yy}/(n-1)$, and the sample covariance is
$S_{xy}/(n-1)$. The correlation from Chapter~4 is
$r = S_{xy}/\sqrt{S_{xx}S_{yy}}$.

\begin{theorem}[Least squares estimates]
\label{thm:slr-ls}
If the $x_i$ are not all equal, $\mathrm{SSE}(b_0,b_1)$ has a unique minimiser
\[
  b_1 = \frac{S_{xy}}{S_{xx}} = r\,\frac{s_y}{s_x},
  \qquad
  b_0 = \ybar - b_1\,\xbar .
\]
\end{theorem}

\begin{proof}
$\mathrm{SSE}$ is a quadratic in $b_0$ and $b_1$ that opens upward, so its
minimum is where both partial derivatives vanish.

\emph{Intercept.} Differentiate with respect to $b_0$:
\[
  \frac{\partial\,\mathrm{SSE}}{\partial b_0}
  = -2\sum_{i=1}^n (y_i - b_0 - b_1 x_i) = 0
  \quad\Longrightarrow\quad
  \sum y_i = n b_0 + b_1 \sum x_i .
\]
Divide by $n$: $\ybar = b_0 + b_1\xbar$. So $b_0 = \ybar - b_1\xbar$.

\emph{Slope.} Differentiate with respect to $b_1$:
\[
  \frac{\partial\,\mathrm{SSE}}{\partial b_1}
  = -2\sum_{i=1}^n x_i (y_i - b_0 - b_1 x_i) = 0 .
\]
Substitute $b_0 = \ybar - b_1\xbar$, so that
$y_i - b_0 - b_1x_i = (y_i-\ybar) - b_1(x_i-\xbar)$:
\[
  \sum_{i=1}^n x_i (y_i-\ybar) = b_1 \sum_{i=1}^n x_i (x_i-\xbar).
\]
Because $\sum (y_i-\ybar) = 0$ and $\sum(x_i-\xbar)=0$, we may replace the
leading $x_i$ by $x_i-\xbar$ on both sides without changing either sum. That
gives $S_{xy} = b_1 S_{xx}$, so $b_1 = S_{xy}/S_{xx}$. We need $S_{xx}>0$,
which holds exactly when the $x_i$ are not all equal.

\emph{The $r$ form.} Since $s_y/s_x = \sqrt{S_{yy}/S_{xx}}$,
\[
  r\,\frac{s_y}{s_x} = \frac{S_{xy}}{\sqrt{S_{xx}S_{yy}}}\cdot\frac{\sqrt{S_{yy}}}{\sqrt{S_{xx}}}
  = \frac{S_{xy}}{S_{xx}} = b_1 .
\]
\end{proof}

The two equations set to zero are called the \textbf{normal equations}. Written
in terms of residuals they say
\[
  \sum_{i=1}^n e_i = 0, \qquad \sum_{i=1}^n x_i e_i = 0 .
\]
Residuals of a least squares fit always sum to zero and are uncorrelated with
$x$. We will use both facts repeatedly.

\begin{recoverybox}
Pause and think: if every $y_i$ is shifted up by $10$, what happens to $b_1$
and $b_0$? (Answer: $S_{xy}$ is unchanged, so $b_1$ is unchanged; $\ybar$ rises
by $10$, so $b_0$ rises by $10$. The line just moves up.)
\end{recoverybox}

\begin{intuitionbox}
The intercept equation $\ybar = b_0 + b_1\xbar$ says the point $(\xbar,\ybar)$
lies on the fitted line. The least squares line always passes through the
\textbf{centroid} of the data. So fitting reduces to one question: through the
centroid, which way should the line tilt? The slope formula
$b_1 = r\,s_y/s_x$ answers it: tilt by $r$ standard deviations of $y$ per
standard deviation of $x$.
\end{intuitionbox}

Because the line pivots on the centroid, $\mathrm{SSE}$ as a function of slope
alone is
$\sum\bigl((y_i-\ybar)-b(x_i-\xbar)\bigr)^2 = S_{yy} - 2bS_{xy} + b^2S_{xx}$,
a parabola with vertex at $b = S_{xy}/S_{xx}$. Figure~\ref{fig:slr-sse-bowl}
plots it for the upload data.

\begin{figure}[H]
\centering
\begin{tikzpicture}
\begin{axis}[
  width=10cm, height=5.8cm, axis lines=left,
  xmin=-0.1, xmax=2.6, ymin=0, ymax=32,
  xlabel={candidate slope $b$ (line through the centroid)},
  ylabel={$\mathrm{SSE}(b)$}, clip=false,
]
\addplot[structurepurple, very thick, domain=-0.05:2.5, samples=80] {28-42*x+17.5*x^2};
\addplot[only marks, mark=*, mark size=2.2pt, black] coordinates {(0,28) (0.6,9.1) (1.2,2.8) (1.8,9.1) (2.4,28)};
\draw[dashed, gray] (axis cs:1.2,0) -- (axis cs:1.2,2.8);
\node[anchor=south, font=\small] (mn) at (axis cs:1.2,17) {minimum: $b_1=1.2$, $\mathrm{SSE}=2.8$};
\draw[-{Stealth}] (mn.south) -- (axis cs:1.2,3.6);
\node[anchor=west, font=\small] at (axis cs:0.16,28.5) {flat line, $\mathrm{SSE}=28=S_{yy}$};
\end{axis}
\end{tikzpicture}
\caption{Sum of squared errors as a function of slope for the upload data,
$\mathrm{SSE}(b) = 28 - 42b + 17.5b^2$. It is a parabola; the bottom is the
least squares slope.}
\label{fig:slr-sse-bowl}
\end{figure}

\section{A worked example by hand}
\label{sec:slr-worked}

\begin{walkbox}
Fit the least squares line to the six uploads.

\textbf{Step 1: means.} $\xbar = 21/6 = 3.5$ MB and $\ybar = 36/6 = 6$ s.

\textbf{Step 2: deviations and products.}

\begin{center}
\begin{tabular}{ccccccc}
\toprule
$x_i$ & $y_i$ & $x_i-\xbar$ & $y_i-\ybar$ & $(x_i-\xbar)^2$ & $(x_i-\xbar)(y_i-\ybar)$ & $(y_i-\ybar)^2$ \\
\midrule
1 & 3 & $-2.5$ & $-3$ & 6.25 & 7.5 & 9 \\
2 & 5 & $-1.5$ & $-1$ & 2.25 & 1.5 & 1 \\
3 & 4 & $-0.5$ & $-2$ & 0.25 & 1.0 & 4 \\
4 & 7 & $0.5$ & $1$ & 0.25 & 0.5 & 1 \\
5 & 8 & $1.5$ & $2$ & 2.25 & 3.0 & 4 \\
6 & 9 & $2.5$ & $3$ & 6.25 & 7.5 & 9 \\
\midrule
 & & 0 & 0 & $S_{xx}=17.5$ & $S_{xy}=21$ & $S_{yy}=28$ \\
\bottomrule
\end{tabular}
\end{center}

\textbf{Step 3: slope and intercept.}
\[
  b_1 = \frac{S_{xy}}{S_{xx}} = \frac{21}{17.5} = 1.2,
  \qquad
  b_0 = \ybar - b_1\xbar = 6 - 1.2\times 3.5 = 1.8 .
\]
The fitted line is $\hat y = 1.8 + 1.2x$.

\textbf{Step 4: check with $r$.} $s_x = \sqrt{17.5/5} = 1.871$,
$s_y=\sqrt{28/5} = 2.366$, $r = 21/\sqrt{17.5\times 28} = 21/\sqrt{490} = 0.9487$.
Then $r\,s_y/s_x = 0.9487\times 2.366/1.871 = 1.2$. Same slope.

\textbf{Step 5: fitted values and residuals.}

\begin{center}
\begin{tabular}{ccccc}
\toprule
$x_i$ & $y_i$ & $\hat y_i = 1.8+1.2x_i$ & $e_i = y_i-\hat y_i$ & $e_i^2$ \\
\midrule
1 & 3 & 3.0 & $0.0$ & 0.00 \\
2 & 5 & 4.2 & $0.8$ & 0.64 \\
3 & 4 & 5.4 & $-1.4$ & 1.96 \\
4 & 7 & 6.6 & $0.4$ & 0.16 \\
5 & 8 & 7.8 & $0.2$ & 0.04 \\
6 & 9 & 9.0 & $0.0$ & 0.00 \\
\midrule
 & & & $\sum e_i = 0$ & $\mathrm{SSE}=2.80$ \\
\bottomrule
\end{tabular}
\end{center}
The residuals sum to zero, as the normal equations promise. So does the
second normal equation:
\[
  \textstyle\sum x_ie_i = 2(0.8) + 3(-1.4) + 4(0.4) + 5(0.2) = 0 .
\]
\end{walkbox}

In practice one function call does all of this.

\begin{lstlisting}[language=Python]
import numpy as np
from scipy import stats
x = np.array([1, 2, 3, 4, 5, 6]); y = np.array([3, 5, 4, 7, 8, 9])
fit = stats.linregress(x, y)
print(fit.slope, fit.intercept)   # 1.2  1.8
print(fit.rvalue**2)              # 0.9  (R^2, Section 12.6)
print(fit.stderr, fit.pvalue)     # 0.2  0.0039 (Section 12.9)
\end{lstlisting}

\section{Interpreting slope and intercept}
\label{sec:slr-interpret}

\textbf{Slope.} $b_1 = 1.2$ s/MB. Each extra megabyte is associated with
$1.2$ more seconds of upload time \emph{on average}. The units of the slope are
always (units of $y$) per (unit of $x$). Change $x$ from MB to KB and the slope
becomes $0.0012$ s/KB. The fit is identical; only the units moved.

\textbf{Intercept.} $b_0 = 1.8$ s is the predicted time at $x=0$: a zero-byte
upload. Here that reading makes sense: it is fixed overhead (connection setup,
handshake). Often it does not. Regress adult weight on height and the intercept
is the predicted weight of a person 0 cm tall, which is meaningless. The
intercept is still needed to place the line; just do not interpret it outside
the range of the data.

\begin{trapbox}
``Associated with'' is not ``causes.'' A regression slope summarises how $y$
and $x$ move together in \emph{these data}. If the data come from an
observational study, a confounder could drive both (Chapter~5). Only a
randomised experiment, where $x$ is set by the experimenter, licenses reading
$b_1$ as ``if I increase $x$ by one, $y$ goes up by $b_1$.'' File size is a
reasonable causal story here, but the regression alone does not prove it.
\end{trapbox}

\section{How good is the fit? \texorpdfstring{$R^2$}{R squared}}
\label{sec:slr-r2}

Before seeing $x$, the best constant prediction of $y$ is $\ybar$, with total
squared error $S_{yy}$. After fitting the line, the squared error drops to
$\mathrm{SSE}$. The fraction of the original squared error that the line
removed measures how useful $x$ is.

\begin{definition}[Sums of squares and $R^2$]
\begin{align*}
  \mathrm{SST} &= \sum (y_i-\ybar)^2 = S_{yy} && \text{(total)},\\
  \mathrm{SSR} &= \sum (\hat y_i-\ybar)^2 && \text{(regression, ``explained'')},\\
  \mathrm{SSE} &= \sum (y_i-\hat y_i)^2 && \text{(error, ``unexplained'')}.
\end{align*}
The \textbf{coefficient of determination} is
$R^2 = \mathrm{SSR}/\mathrm{SST} = 1 - \mathrm{SSE}/\mathrm{SST}$.
\end{definition}

\begin{theorem}[Sum of squares decomposition]
For a least squares fit with an intercept,
$\mathrm{SST} = \mathrm{SSR} + \mathrm{SSE}$, and $R^2 = r^2$.
\end{theorem}

\begin{proof}
Split each deviation: $y_i - \ybar = (\hat y_i - \ybar) + e_i$. Square and sum:
\[
  \mathrm{SST} = \mathrm{SSR} + \mathrm{SSE} + 2\sum_{i=1}^n (\hat y_i-\ybar)e_i .
\]
Since $\hat y_i - \ybar = b_1(x_i-\xbar)$, the cross term is
$2b_1\bigl(\sum x_ie_i - \xbar\sum e_i\bigr)$, and both sums vanish by the
normal equations. For the second claim,
$\mathrm{SSR} = b_1^2 S_{xx} = S_{xy}^2/S_{xx}$, so
\[
  R^2 = \frac{S_{xy}^2}{S_{xx}S_{yy}} = r^2. \qedhere
\]
\end{proof}

Figure~\ref{fig:slr-decomp} shows the split for a single point. The cross term
vanishing is a Pythagorean theorem: the vector of residuals is perpendicular to
the vector of fitted deviations.

\begin{figure}[H]
\centering
\begin{tikzpicture}
\begin{axis}[
  width=11cm, height=6.5cm, axis lines=left,
  xmin=0, xmax=7.6, ymin=1, ymax=10.5,
  xlabel={file size $x$ (MB)}, ylabel={upload time $y$ (s)}, clip=false,
]
\addplot[sourceblue, very thick, domain=0.3:6.7] {1.8+1.2*x};
\addplot[gray, dashed, thick, domain=0.3:6.7] {6};
\node[gray, anchor=west, font=\small] at (axis cs:6.7,6) {$\ybar=6$};
\node[sourceblue, anchor=west, font=\small] at (axis cs:6.7,9.84) {$\hat y$};
\addplot[only marks, mark=*, mark size=1.8pt, gray!60] coordinates {(1,3) (2,5) (3,4) (4,7) (6,9)};
\addplot[only marks, mark=*, mark size=2.6pt, black] coordinates {(5,8)};
\draw[decorate, decoration={brace, amplitude=5pt, mirror}, targetgreen, thick] (axis cs:5.12,6) -- (axis cs:5.12,7.8);
\draw[decorate, decoration={brace, amplitude=3pt, mirror}, trapred, thick] (axis cs:5.12,7.82) -- (axis cs:5.12,8.0);
\draw[decorate, decoration={brace, amplitude=6pt}, structurepurple, thick] (axis cs:4.85,6) -- (axis cs:4.85,8);
\node[targetgreen, anchor=west, font=\small] (ex) at (axis cs:5.6,4.6) {$\hat y_5-\ybar = 1.8$ (explained)};
\draw[targetgreen, -{Stealth}] (ex.north west) -- (axis cs:5.32,6.8);
\node[trapred, anchor=east, font=\small] (re) at (axis cs:3.9,9.6) {$y_5-\hat y_5 = 0.2$ (residual)};
\draw[trapred, -{Stealth}] (re.east) -- (axis cs:5.2,7.95);
\node[structurepurple, anchor=east, font=\small] (to) at (axis cs:4.4,2.4) {$y_5-\ybar = 2$ (total)};
\draw[structurepurple, -{Stealth}] (to.north east) -- (axis cs:4.68,6.9);
\end{axis}
\end{tikzpicture}
\caption{For point $(5,8)$ the total deviation from the mean ($2$) splits into
the part the line explains ($1.8$) and the residual ($0.2$). Summing the
\emph{squares} of these pieces over all points gives
$\mathrm{SST}=\mathrm{SSR}+\mathrm{SSE}$.}
\label{fig:slr-decomp}
\end{figure}

For the uploads: $\mathrm{SST}=28$, $\mathrm{SSE}=2.8$, so
$\mathrm{SSR}=25.2$ and
\[
  R^2 = \frac{25.2}{28} = 0.90 = (0.9487)^2 = r^2 .
\]
File size explains $90\%$ of the variation in upload time. The other $10\%$ is
network noise and anything else we did not record.

\begin{recoverybox}
Pause and think: what is $R^2$ if the fitted slope is exactly $0$? And if every
point lies on the line? (Answers: $0$, because then $\hat y_i = \ybar$ and
$\mathrm{SSR}=0$; and $1$, because $\mathrm{SSE}=0$.)
\end{recoverybox}

\begin{figure}[H]
\centering
\begin{tikzpicture}
\begin{axis}[
  name=left, width=6.3cm, height=5.6cm, axis lines=left,
  xmin=0, xmax=10.8, ymin=-3, ymax=18, xlabel={$x$}, ylabel={$y$},
  title style={font=\small}, title={\small $R^2 = 0.92$},
]
\addplot[sourceblue, very thick, domain=0.5:10.5] {1.357+1.120*x};
\addplot[only marks, mark=*, mark size=1.5pt] coordinates {(1.00,0.81) (1.47,2.83) (1.95,4.13) (2.42,3.71) (2.89,4.57) (3.37,5.22) (3.84,5.89) (4.32,7.34) (4.79,6.76) (5.26,7.00) (5.74,8.28) (6.21,10.29) (6.68,8.34) (7.16,9.67) (7.63,8.05) (8.11,11.45) (8.58,11.24) (9.05,12.39) (9.53,11.47) (10.00,10.96)};
\end{axis}
\begin{axis}[
  at={(left.east)}, anchor=west, xshift=0.5cm,
  width=6.3cm, height=5.6cm, axis lines=left,
  xmin=0, xmax=10.8, ymin=-3, ymax=18, xlabel={$x$},
  title style={font=\small}, title={\small $R^2 = 0.26$},
]
\addplot[sourceblue, very thick, domain=0.5:10.5] {2.401+0.895*x};
\addplot[only marks, mark=*, mark size=1.5pt] coordinates {(1.00,6.54) (1.47,6.49) (1.95,0.48) (2.42,-1.47) (2.89,6.69) (3.37,4.68) (3.84,16.29) (4.32,3.86) (4.79,4.68) (5.26,6.82) (5.74,1.08) (6.21,4.65) (6.68,15.26) (7.16,8.63) (7.63,7.20) (8.11,5.02) (8.58,12.90) (9.05,13.69) (9.53,12.55) (10.00,10.44)};
\end{axis}
\end{tikzpicture}
\caption{Similar slopes, very different $R^2$. On the right the trend is real
but most of the variation in $y$ comes from noise that $x$ does not explain.}
\label{fig:slr-r2}
\end{figure}

\begin{trapbox}
$R^2$ is not a grade for the model. A low $R^2$ can come with a precisely
estimated, important slope (a drug that lowers blood pressure by 5 mmHg on
average, in a population where blood pressure varies by 20). A high $R^2$ can
come with a wrong model (a straight line through a curve, Section~\ref{sec:slr-diag}).
$R^2$ also says nothing about causation, and it depends on the range of $x$:
sample a wider range of $x$ and $R^2$ goes up even though $\sigma$ is unchanged.
\end{trapbox}

\section{Regression to the mean}
\label{sec:slr-rtm}

Write the fitted line in standard units. Let
$z_x = (x-\xbar)/s_x$ and $z_y = (y-\ybar)/s_y$. Using $b_1 = r\,s_y/s_x$ and
the fact that the line passes through the centroid,
\[
  \hat z_y = r\, z_x .
\]
Since $\abs{r}\le 1$, the prediction is always \emph{closer to the mean} (in SD
units) than the input. A point $2$ SDs above average in $x$ is predicted to be
only $2r$ SDs above average in $y$.

\begin{example}[Parents and children]
Fifty parent--adult child pairs (simulated, heights in cm) have
$\xbar = 170.0$, $s_x = 6.16$ for parents, $\ybar=170.6$, $s_y = 6.09$ for
children, and $r = 0.34$. The slope is $b_1 = r\,s_y/s_x = 0.339$, so the
fitted line is $\hat y = 113.0 + 0.339x$. A parent at $182$ cm is
$(182-170.0)/6.16 = 1.95$ SDs above average. Predicted child height:
$170.6 + 0.339\times(182-170.0) \approx 174.6$ cm, only
$0.34\times 1.95 = 0.67$ SDs above the children's average.
\end{example}

\begin{historybox}
In 1805 Adrien-Marie Legendre was trying to pin down the orbits of comets from
more observations than unknowns: the equations could not all hold at once
because every measurement had error. He proposed choosing the answer that
makes the sum of squared mismatches smallest, and called it the ``method of
least squares.'' Carl Friedrich Gauss published it in 1809, claimed he had
used it since 1795, and showed it is the best choice when errors follow the
bell curve that now carries his name (Section~\ref{sec:slr-mle}). The word
``regression'' came eighty years later. Francis Galton measured the heights of
928 adult children and their parents and found in 1886 that children of
unusually tall or short parents were, on average, closer to the middle. He
called it ``regression towards mediocrity'' (``mediocrity'' meaning the
average). The name stuck to the whole method of fitting lines, even though the
effect it describes is a property of imperfect correlation, not of lines.
\end{historybox}

Nothing biological pulls children toward the average. Tall parents are tall partly because of genes they pass on and partly
because of luck that they do not. The luck does not repeat, so on average the
children are less extreme.

\begin{figure}[H]
\centering
\begin{tikzpicture}
\begin{axis}[
  width=10cm, height=9cm, axis lines=left,
  xmin=152, xmax=190, ymin=152, ymax=190,
  xlabel={parent height $x$ (cm)}, ylabel={child height $y$ (cm)},
  legend pos=north west, legend style={font=\small, cells={anchor=west}},
]
\addplot[only marks, mark=*, mark size=1.4pt, gray] coordinates {
(170.0,167.7) (177.3,176.3) (175.2,178.5) (175.1,169.2) (181.3,183.1) (161.6,175.6) (165.6,174.3) (160.8,172.4) (169.2,176.3) (177.0,187.2) (169.8,171.0) (173.5,171.7) (156.6,167.0) (171.0,165.0) (163.7,157.4) (182.4,170.9) (176.2,175.4) (176.6,176.2) (169.6,160.5) (174.3,160.7) (174.6,170.4) (167.6,167.7) (166.5,168.0) (169.2,173.7) (165.8,175.3) (165.8,169.3) (168.0,172.7) (164.9,160.8) (175.4,165.6) (158.8,165.8) (175.8,176.9) (165.6,172.2) (166.2,171.7) (160.5,170.0) (169.0,174.4) (168.3,158.5) (171.3,166.2) (166.3,170.9) (170.7,173.9) (182.7,167.2) (172.9,178.5) (166.0,165.3) (176.7,171.9) (169.1,171.2) (174.2,176.6) (174.3,163.5) (167.3,174.6) (156.5,163.7) (167.6,164.1) (173.9,170.3)};
\addlegendentry{data}
\addplot[sourceblue, very thick, domain=152:190] {113.011+0.33851*x};
\addlegendentry{$y$ on $x$ (slope $0.34$)}
\addplot[trapred, very thick, domain=165.4:175.7] {-319.16+2.8812*x};
\addlegendentry{$x$ on $y$ (slope $1/0.35$)}
\addplot[black, dashed, thick, domain=152:190] {170.552+0.98768*(x-169.961)};
\addlegendentry{SD line (slope $s_y/s_x$)}
\addplot[only marks, mark=square*, mark size=2.5pt, quotientorange] coordinates {(169.96,170.55)};
\end{axis}
\end{tikzpicture}
\caption{Two regression lines for the same data. Blue predicts child from
parent; red predicts parent from child (redrawn in these axes). Both pass
through the centroid (orange) and both are flatter, \emph{in their own
direction}, than the dashed SD line. That flattening is regression to the mean.}
\label{fig:slr-two-lines}
\end{figure}

\paragraph{Why $y$ on $x$ differs from $x$ on $y$.}
Regressing $x$ on $y$ minimises \emph{horizontal} residuals and gives
$x = c_0 + c_1 y$ with $c_1 = r\,s_x/s_y$. If you solve this for $y$, its slope
in the $(x,y)$ plane is $1/c_1 = s_y/(r\,s_x)$. The $y$-on-$x$ slope is
$r\,s_y/s_x$. These agree only when $r=\pm1$. For $\abs{r}<1$ the two lines
sit on either side of the SD line (Figure~\ref{fig:slr-two-lines}). For the
uploads, $x$ on $y$ gives $\hat x = -1 + 0.75y$, i.e.\ $y = 1.33 + 1.33x$ when
rearranged, steeper than $1.8+1.2x$.

\begin{trapbox}
Do not invert a regression line to predict $x$ from $y$. If you want to predict
parent height from child height, fit the regression of $x$ on $y$. Each line is
the best predictor in one direction only.
\end{trapbox}

\begin{happeningbox}
Chapter~11 showed the practical cost: whenever units are selected for being
extreme (the worst-performing servers this week, the lowest pretest scores),
they drift back toward average with no intervention at all, and a slope of
$r<1$ is the reason. The fix is a control group selected the same way
(Chapter~5).
\end{happeningbox}

\section{Residual plots and diagnostics}
\label{sec:slr-diag}

$b_0$, $b_1$ and $R^2$ can be computed for any cloud of points. Whether a
straight line is a sensible summary is a separate question, and the residuals
answer it. The basic tool is the \textbf{residual plot}: residuals $e_i$ on the
vertical axis against fitted values $\hat y_i$ (or against $x_i$; for one
predictor $\hat y_i$ is a linear function of $x_i$, so the two plots have the
same pattern).

If the model is right, the residuals are just noise: a horizontal band centred
on $0$, with no pattern and constant width.

\begin{figure}[H]
\centering
\begin{tikzpicture}
\begin{axis}[
  name=left, width=6.3cm, height=5.2cm, axis lines=left,
  xmin=0, xmax=10.5, ymin=0, ymax=15, xlabel={$x$}, ylabel={$y$},
  title style={font=\small}, title={\small data and fitted line},
]
\addplot[sourceblue, very thick, domain=0.3:10.3] {2.798+1.0074*x};
\addplot[only marks, mark=*, mark size=1.5pt] coordinates {(0.50,3.81) (0.90,4.28) (1.29,3.61) (1.69,5.31) (2.08,2.97) (2.48,4.19) (2.88,5.57) (3.27,6.16) (3.67,5.50) (4.06,7.39) (4.46,8.51) (4.85,8.31) (5.25,8.47) (5.65,9.47) (6.04,9.45) (6.44,8.30) (6.83,9.63) (7.23,11.52) (7.62,9.49) (8.02,11.08) (8.42,11.03) (8.81,10.91) (9.21,11.20) (9.60,12.71) (10.00,13.32)};
\end{axis}
\begin{axis}[
  at={(left.east)}, anchor=west, xshift=0.7cm,
  width=6.3cm, height=5.2cm, axis lines=left,
  xmin=2.5, xmax=13.5, ymin=-3, ymax=3, xlabel={fitted $\hat y$}, ylabel={residual $e$},
  title style={font=\small}, title={\small residual plot: healthy},
]
\addplot[gray, dashed, domain=2.5:13.5] {0};
\addplot[only marks, mark=*, mark size=1.5pt, targetgreen] coordinates {(3.30,0.50) (3.70,0.58) (4.10,-0.49) (4.50,0.81) (4.90,-1.93) (5.30,-1.11) (5.69,-0.12) (6.09,0.07) (6.49,-0.99) (6.89,0.50) (7.29,1.22) (7.69,0.62) (8.09,0.39) (8.49,0.99) (8.88,0.56) (9.28,-0.99) (9.68,-0.05) (10.08,1.44) (10.48,-0.99) (10.88,0.20) (11.28,-0.25) (11.68,-0.76) (12.07,-0.88) (12.47,0.23) (12.87,0.44)};
\end{axis}
\end{tikzpicture}
\caption{A model that fits. The residual plot is a patternless horizontal band
around zero with roughly constant width.}
\label{fig:slr-resid-good}
\end{figure}

Three patterns signal trouble.

\textbf{Curvature.} A U or arch in the residuals means the true relationship is
not a line. The line is too high in the middle and too low at the ends (or the
reverse). Figure~\ref{fig:slr-resid-curve} shows a quadratic relationship fit
by a line. The fit has $R^2 = 0.92$, which looks excellent, yet the residuals
reveal the model is wrong. Fix: transform $x$ or $y$, or add an $x^2$ term
(Chapter~13).

\begin{figure}[H]
\centering
\begin{tikzpicture}
\begin{axis}[
  name=left, width=6.3cm, height=5.2cm, axis lines=left,
  xmin=0, xmax=10.5, ymin=-5, ymax=28, xlabel={$x$}, ylabel={$y$},
  title style={font=\small}, title={\small line through a curve},
]
\addplot[sourceblue, very thick, domain=0.3:10.3] {-3.939+2.6382*x};
\addplot[only marks, mark=*, mark size=1.5pt] coordinates {(0.50,1.31) (0.90,0.04) (1.29,1.19) (1.69,1.99) (2.08,2.84) (2.48,3.07) (2.88,3.73) (3.27,4.55) (3.67,5.02) (4.06,4.47) (4.46,6.02) (4.85,6.32) (5.25,5.99) (5.65,7.62) (6.04,11.80) (6.44,10.47) (6.83,10.88) (7.23,13.60) (7.62,16.75) (8.02,17.05) (8.42,19.04) (8.81,20.20) (9.21,22.13) (9.60,25.07) (10.00,26.64)};
\end{axis}
\begin{axis}[
  at={(left.east)}, anchor=west, xshift=0.7cm,
  width=6.3cm, height=5.2cm, axis lines=left,
  xmin=-3.5, xmax=23.5, ymin=-5, ymax=5, xlabel={fitted $\hat y$}, ylabel={residual $e$},
  title style={font=\small}, title={\small residuals: U shape},
]
\addplot[gray, dashed, domain=-3.5:23.5] {0};
\addplot[only marks, mark=*, mark size=1.5pt, trapred] coordinates {(-2.62,3.93) (-1.58,1.62) (-0.53,1.72) (0.51,1.48) (1.56,1.29) (2.60,0.47) (3.65,0.09) (4.69,-0.14) (5.73,-0.71) (6.78,-2.30) (7.82,-1.81) (8.87,-2.55) (9.91,-3.93) (10.96,-3.33) (12.00,-0.20) (13.04,-2.58) (14.09,-3.21) (15.13,-1.54) (16.18,0.57) (17.22,-0.17) (18.27,0.77) (19.31,0.89) (20.35,1.78) (21.40,3.67) (22.44,4.19)};
\end{axis}
\end{tikzpicture}
\caption{Curvature. The scatter plot hides it (the line ``looks fine'' and
$R^2=0.92$), but the residual plot shows a clear U: positive, then negative,
then positive.}
\label{fig:slr-resid-curve}
\end{figure}

\textbf{Fan shape (heteroscedasticity).} Residuals spread out as $\hat y$
grows. The noise is not constant: big files have more variable upload times
than small ones. The line itself may be fine, but $\SE(b_1)$ and every interval
in Section~\ref{sec:slr-inference} assume constant $\sigma$ and will be wrong.
Fix: model $\log y$, use weighted least squares, or use robust
(``sandwich'') standard errors.

\begin{figure}[H]
\centering
\begin{tikzpicture}
\begin{axis}[
  name=left, width=6.3cm, height=5.2cm, axis lines=left,
  xmin=0, xmax=10.5, ymin=0, ymax=16, xlabel={$x$}, ylabel={$y$},
  title style={font=\small}, title={\small spread grows with $x$},
]
\addplot[sourceblue, very thick, domain=0.3:10.3] {2.874+1.0132*x};
\addplot[only marks, mark=*, mark size=1.5pt] coordinates {(0.50,3.49) (0.90,4.10) (1.29,4.64) (1.69,4.60) (2.08,5.22) (2.48,5.64) (2.88,5.39) (3.27,7.21) (3.67,6.51) (4.06,6.43) (4.46,5.59) (4.85,6.32) (5.25,10.25) (5.65,7.04) (6.04,9.69) (6.44,8.85) (6.83,10.71) (7.23,11.32) (7.62,9.15) (8.02,10.44) (8.42,11.63) (8.81,12.06) (9.21,12.14) (9.60,11.85) (10.00,14.57)};
\end{axis}
\begin{axis}[
  at={(left.east)}, anchor=west, xshift=0.7cm,
  width=6.3cm, height=5.2cm, axis lines=left,
  xmin=2.5, xmax=13.5, ymin=-3, ymax=3, xlabel={fitted $\hat y$}, ylabel={residual $e$},
  title style={font=\small}, title={\small residuals: fan},
]
\addplot[gray, dashed, domain=2.5:13.5] {0};
\addplot[quotientorange, dashed, domain=3:13.3] {0.2*x};
\addplot[quotientorange, dashed, domain=3:13.3] {-0.2*x};
\addplot[only marks, mark=*, mark size=1.5pt, trapred] coordinates {(3.38,0.11) (3.78,0.31) (4.18,0.46) (4.58,0.01) (4.98,0.24) (5.39,0.25) (5.79,-0.39) (6.19,1.02) (6.59,-0.08) (6.99,-0.56) (7.39,-1.80) (7.79,-1.47) (8.19,2.06) (8.59,-1.56) (9.00,0.70) (9.40,-0.55) (9.80,0.91) (10.20,1.12) (10.60,-1.45) (11.00,-0.57) (11.40,0.22) (11.80,0.26) (12.20,-0.06) (12.60,-0.75) (13.01,1.57)};
\end{axis}
\end{tikzpicture}
\caption{Fan shape. Residuals are tight for small fitted values and wide for
large ones (orange guides). The constant-variance assumption fails.}
\label{fig:slr-resid-fan}
\end{figure}

\textbf{Outliers, leverage, and influence.} These are three different things.

\begin{definition}[Outlier, leverage, influence]
An \textbf{outlier} is a point with a large residual: its $y$ is far from the
line. A point has high \textbf{leverage} if its $x$ is far from $\xbar$; the
leverage of point $i$ is
\[
  h_i = \frac{1}{n} + \frac{(x_i-\xbar)^2}{S_{xx}} .
\]
A point is \textbf{influential} if deleting it changes the fitted line a lot.
Influence needs both: high leverage \emph{and} a $y$ value off the line the
other points suggest.
\end{definition}

Figure~\ref{fig:slr-influence} adds one extra point to ten well-behaved points
(base fit $\hat y = 1.49 + 0.741x$) in three ways.

\begin{figure}[H]
\centering
\begin{tikzpicture}
\pgfplotsset{inflpanel/.style={width=4.75cm, height=4.6cm, axis lines=left,
  xmin=0, xmax=21, ymin=0, ymax=18, xtick={0,10,20}, ytick={0,5,10,15},
  title style={font=\small, align=center}, tick label style={font=\footnotesize}}}
\begin{axis}[inflpanel, title style={font=\small}, name=p1, title={A: outlier,\\ low leverage}, ylabel={$y$}, xlabel={$x$}]
\addplot[gray, thick, dashed, domain=0:21] {1.4925+0.74066*x};
\addplot[trapred, thick, domain=0:21] {2.0774+0.74066*x};
\addplot[only marks, mark=*, mark size=1.3pt] coordinates {(1.00,2.59) (2.00,2.92) (3.00,2.96) (4.00,4.44) (5.00,5.04) (6.00,6.08) (7.00,7.71) (8.00,7.27) (9.00,8.19) (10.00,8.47)};
\addplot[only marks, mark=triangle*, mark size=3pt, trapred] coordinates {(5.5,12)};
\end{axis}
\begin{axis}[inflpanel, title style={font=\small}, name=p2, at={(p1.east)}, anchor=west, xshift=0.3cm, yticklabels={}, title={B: high leverage,\\ on trend}, xlabel={$x$}]
\addplot[gray, thick, dashed, domain=0:21] {1.4925+0.74066*x};
\addplot[trapred, thick, domain=0:21] {1.3276+0.77411*x};
\addplot[only marks, mark=*, mark size=1.3pt] coordinates {(1.00,2.59) (2.00,2.92) (3.00,2.96) (4.00,4.44) (5.00,5.04) (6.00,6.08) (7.00,7.71) (8.00,7.27) (9.00,8.19) (10.00,8.47)};
\addplot[only marks, mark=triangle*, mark size=3pt, trapred] coordinates {(20,17)};
\end{axis}
\begin{axis}[inflpanel, title style={font=\small}, name=p3, at={(p2.east)}, anchor=west, xshift=0.3cm, yticklabels={}, title={C: high leverage,\\ off trend}, xlabel={$x$}]
\addplot[gray, thick, dashed, domain=0:21] {1.4925+0.74066*x};
\addplot[trapred, thick, domain=0:21] {4.4156+0.14786*x};
\addplot[only marks, mark=*, mark size=1.3pt] coordinates {(1.00,2.59) (2.00,2.92) (3.00,2.96) (4.00,4.44) (5.00,5.04) (6.00,6.08) (7.00,7.71) (8.00,7.27) (9.00,8.19) (10.00,8.47)};
\addplot[only marks, mark=triangle*, mark size=3pt, trapred] coordinates {(20,4)};
\end{axis}
\end{tikzpicture}
\caption{Dashed grey: fit to the ten black points. Red: fit after adding the
red triangle. A (outlier near $\xbar$) only lifts the line; slope stays
$0.741$. B (far in $x$ but on trend) barely moves it; slope $0.774$. C (far in
$x$ and off trend) drags the slope from $0.741$ to $0.148$. Only C is
influential.}
\label{fig:slr-influence}
\end{figure}

\begin{whybox}
Why does a far-out $x$ give a point so much pull? The slope is a weighted
average of the points' individual slopes from the centroid:
\[
  b_1 = \sum_{i=1}^n \frac{(x_i-\xbar)^2}{S_{xx}}\cdot\frac{y_i-\ybar}{x_i-\xbar},
\]
with weight proportional to $(x_i-\xbar)^2$ (a point with $x_i=\xbar$ gets
weight $0$ and is left out of the sum). A point ten units from $\xbar$
gets a hundred times the weight of a point one unit away. It acts like a long
lever arm on a seesaw pivoting at the centroid.
\end{whybox}

A standard numerical summary of influence is \textbf{Cook's distance}, which
measures how much all fitted values move when point $i$ is deleted; values
near or above $1$ deserve a look. Never delete a point just because it is
influential. Find out why it is different: a data-entry error, a different
population, or the most informative observation you have.

\section{Inference for the slope}
\label{sec:slr-inference}

The slope $b_1 = 1.2$ came from six uploads. Six different uploads would give a
different slope. Just as $\xbar$ has a sampling distribution (Chapter~6), so
does $b_1$. Figure~\ref{fig:slr-many-lines} simulates it: twenty samples of six
uploads at the same file sizes, from a true line $y = 2 + 1.2x$ with
$\sigma = 0.8$.

\begin{figure}[H]
\centering
\begin{tikzpicture}
\begin{axis}[
  width=10.5cm, height=6.5cm, axis lines=left,
  xmin=0, xmax=7, ymin=0, ymax=13,
  xlabel={file size $x$ (MB)}, ylabel={upload time $y$ (s)},
]
\foreach \a/\b in {1.537/1.424, 1.923/1.250, 1.725/1.171, 0.654/1.564, 0.469/1.623, 4.088/0.689, 3.028/0.994, 3.251/0.865, 1.727/1.233, 2.769/1.127, 1.567/1.352, 1.413/1.371, 3.135/0.870, 2.628/1.179, 3.129/0.894, 0.623/1.569, 2.818/0.943, 2.428/1.180, 2.457/0.963, 1.070/1.491} {
  \addplot[sourceblue, opacity=0.45, domain=0.5:6.5] {\a+\b*x};
}
\addplot[black, ultra thick, domain=0.5:6.5] {2+1.2*x};
\addplot[only marks, mark=square*, mark size=2.3pt, quotientorange] coordinates {(3.5,6.2)};
\node[anchor=north west, font=\small] at (axis cs:0.6,12.5) {black: true line $y=2+1.2x$};
\node[anchor=north west, font=\small, sourceblue] at (axis cs:0.6,11.4) {blue: 20 fitted lines};
\end{axis}
\end{tikzpicture}
\caption{Twenty least squares lines, each fitted to a fresh sample of six
points from the same model. The lines wobble around the truth. They pinch
together near $\xbar=3.5$ (orange) and fan out at the ends, because a slope
error matters more far from the pivot.}
\label{fig:slr-many-lines}
\end{figure}

\begin{theorem}[Sampling distribution of $b_1$]
\label{thm:slr-b1}
Treat the $x_i$ as fixed and assume the model of Section~\ref{sec:slr-model}.
Then
\[
  \E[b_1] = \beta_1, \qquad \SD(b_1) = \frac{\sigma}{\sqrt{S_{xx}}} .
\]
If the errors are Normal, $b_1$ is exactly Normal.
\end{theorem}

\begin{proof}
Because $\sum(x_i-\xbar)\ybar = 0$, we can write
$b_1 = \sum_{i} w_i y_i$ with weights $w_i = (x_i-\xbar)/S_{xx}$. These weights
satisfy $\sum w_i = 0$, $\sum w_i x_i = 1$, and $\sum w_i^2 = 1/S_{xx}$.
Substituting $y_i = \beta_0+\beta_1x_i+\eps_i$,
\[
  b_1 = \beta_0\sum w_i + \beta_1 \sum w_i x_i + \sum w_i\eps_i
      = \beta_1 + \sum w_i \eps_i .
\]
The noise term has mean $0$, so $\E[b_1]=\beta_1$. Independent errors with
variance $\sigma^2$ give $\Var(b_1) = \sigma^2\sum w_i^2 = \sigma^2/S_{xx}$. A
linear combination of independent Normals is Normal.
\end{proof}

The formula $\sigma/\sqrt{S_{xx}}$ tells you how to design an experiment.
The slope is pinned down better when the noise $\sigma$ is small, when there
are more points (each adds to $S_{xx}$), and when the $x$ values are spread out.
Testing upload time only on files between 3.4 and 3.6 MB would tell you almost
nothing about the slope. Note $S_{xx} = (n-1)s_x^2$, so
$\SD(b_1) = \sigma/(s_x\sqrt{n-1})$: the familiar $1/\sqrt n$ again.

\begin{figure}[H]
\centering
\begin{tikzpicture}
\begin{axis}[
  width=10cm, height=5.6cm, axis lines=left,
  xmin=0.35, xmax=2.05, ymin=0, ymax=2.8,
  xlabel={fitted slope $b_1$}, ylabel={density},
]
\addplot[ybar interval, fill=histfill, draw=white] coordinates {(0.40,0.000) (0.50,0.015) (0.60,0.055) (0.70,0.150) (0.80,0.415) (0.90,0.880) (1.00,1.411) (1.10,2.146) (1.20,2.121) (1.30,1.531) (1.40,0.780) (1.50,0.360) (1.60,0.115) (1.70,0.015) (1.80,0.005) (1.90,0.000) (2.00,0.000)};
\addplot[densityred, very thick, domain=0.4:2.0, samples=100] {1/(0.19124*2.50663)*exp(-(x-1.2)^2/(2*0.036573))};
\draw[dashed] (axis cs:1.2,0) -- (axis cs:1.2,2.45);
\node[anchor=south, font=\small] at (axis cs:1.2,2.45) {$\beta_1=1.2$};
\node[densityred, anchor=west, font=\small] at (axis cs:1.45,1.6) {$\Normal\bigl(1.2,\ 0.191^2\bigr)$};
\end{axis}
\end{tikzpicture}
\caption{Histogram of $b_1$ over $2000$ simulated samples of six uploads. It
matches the Normal curve with mean $\beta_1=1.2$ and SD
$\sigma/\sqrt{S_{xx}} = 0.8/\sqrt{17.5} = 0.191$ predicted by
Theorem~\ref{thm:slr-b1}.}
\label{fig:slr-b1-hist}
\end{figure}

\subsection{Estimating \texorpdfstring{$\sigma$}{sigma} and the standard error}

We do not know $\sigma$. Estimate it from the residuals:
\[
  s = \sqrt{\frac{\mathrm{SSE}}{n-2}} = \sqrt{\frac{\sum e_i^2}{n-2}} .
\]
This is called the \textbf{residual standard error}. We divide by $n-2$, not
$n$, for the same reason the sample variance divides by $n-1$ (Chapter~2): the
residuals are measured from a line fitted to the same data, so they are too
small on average. Fitting used up two degrees of freedom (one each for $b_0$
and $b_1$); the residuals satisfy two linear constraints (the normal
equations), leaving $n-2$ free. With this divisor, $\E[s^2] = \sigma^2$.

Plugging in gives the standard error:
\[
  \SE(b_1) = \frac{s}{\sqrt{S_{xx}}} .
\]

\begin{theorem}[$t$ statistic for the slope]
Under the Normal-error model,
\[
  T = \frac{b_1 - \beta_1}{\SE(b_1)} \sim t_{n-2} .
\]
\end{theorem}

This is the same move as the one-sample $t$ test of Chapter~9: replace the
unknown $\sigma$ by $s$ and pay for it with a $t$ distribution, here with $n-2$
degrees of freedom. A $100(1-\alpha)\%$ confidence interval for the slope is
\[
  b_1 \pm t^*_{n-2}\,\SE(b_1),
\]
where $t^*_{n-2}$ is the $t_{n-2}$ critical value ($t_{\alpha/2,\,n-2}$ in the notation of Chapter~7),
and the test of $H_0\colon \beta_1 = 0$ (``$x$ has no linear association with
$y$'') uses $t = b_1/\SE(b_1)$.

\begin{walkbox}
Is upload time really related to file size, and how precisely do we know the
slope?

\textbf{Residual SE.} $s = \sqrt{2.8/4} = \sqrt{0.7} = 0.837$ s.

\textbf{SE of slope.} $\SE(b_1) = 0.837/\sqrt{17.5} = 0.837/4.183 = 0.200$ s/MB.

\textbf{Test.} $t = 1.2/0.2 = 6.0$ on $n-2 = 4$ degrees of freedom. The
two-sided $p$-value is $2\P(t_4 > 6.0) = 0.0039$. Reject $H_0\colon\beta_1=0$
at any conventional level.

\textbf{95\% CI.} $t^*_4 = 2.776$. The interval is
$1.2 \pm 2.776\times 0.2 = 1.2\pm 0.555 = [0.645,\ 1.755]$ s/MB.

\textbf{Reading.} Each extra MB adds between about $0.65$ and $1.76$ seconds on
average. With only six points the interval is wide, but it excludes $0$.
\end{walkbox}

\begin{figure}[H]
\centering
\begin{tikzpicture}
\begin{axis}[
  width=11cm, height=5.4cm, axis lines=left,
  xmin=-7, xmax=7, ymin=0, ymax=0.45,
  xlabel={$t$}, ylabel={density of $t_4$}, clip=false,
  xtick={-6,-2.776,0,2.776,6}, xticklabels={$-6$,$-2.78$,$0$,$2.78$,$6$},
]
\addplot[name path=f, distblue, very thick, domain=-7:7, samples=200] {0.375*(1+x^2/4)^(-2.5)};
\addplot[draw=none, fill=trapred, fill opacity=0.35, domain=2.776:7, samples=60] {0.375*(1+x^2/4)^(-2.5)} \closedcycle;
\addplot[draw=none, fill=trapred, fill opacity=0.35, domain=-7:-2.776, samples=60] {0.375*(1+x^2/4)^(-2.5)} \closedcycle;
\draw[targetgreen, very thick] (axis cs:6,0) -- (axis cs:6,0.12);
\node[targetgreen, anchor=south, font=\small, align=center] at (axis cs:6,0.12) {observed\\$t=6.0$};
\node[trapred, anchor=south, font=\small] at (axis cs:-4.6,0.05) {2.5\%};
\node[trapred, anchor=south, font=\small] at (axis cs:4.0,0.05) {2.5\%};
\end{axis}
\end{tikzpicture}
\caption{The $t_4$ distribution under $H_0\colon\beta_1=0$. The red tails
(beyond $\pm2.776$) hold $5\%$ of the probability. The observed $t=6$ lies far
in the tail; the area beyond $\pm 6$ is $p=0.0039$.}
\label{fig:slr-t-test}
\end{figure}

\begin{trapbox}
Testing $\beta_1 = 0$ is the same as testing $\rho = 0$ (zero population
correlation): $t = b_1/\SE(b_1) = r\sqrt{n-2}/\sqrt{1-r^2}$. Check:
$0.9487\times 2/\sqrt{0.1} = 6.0$. A tiny $p$-value says the slope is not zero.
It does not say the slope is large, that the line fits well, or that $x$
causes $y$.
\end{trapbox}

\section{Predicting: mean response vs.\ a new observation}
\label{sec:slr-predict}

Suppose a 4 MB file is about to be uploaded. The point prediction is
$\hat y_0 = 1.8 + 1.2\times 4 = 6.6$ s. Two different questions have this
same point estimate but different error bars.

\begin{enumerate}
  \item \textbf{Mean response.} What is the \emph{average} upload time of all
  4 MB files, $\beta_0+\beta_1\cdot 4$? Uncertainty comes only from not
  knowing the true line.
  \item \textbf{Prediction.} How long will \emph{this one} 4 MB upload take?
  Now add the noise $\eps$ of the single new upload on top.
\end{enumerate}

Write $\hat y_0 = \ybar + b_1(x_0-\xbar)$. The two pieces $\ybar$ and $b_1$ are
uncorrelated (because $\sum w_i = 0$), so
\[
  \Var(\hat y_0) = \frac{\sigma^2}{n} + (x_0-\xbar)^2\frac{\sigma^2}{S_{xx}} .
\]
A new observation is $y_0 = \beta_0+\beta_1x_0+\eps_0$ with fresh noise
independent of the data, so $\Var(y_0 - \hat y_0) = \sigma^2 + \Var(\hat y_0)$.
Replacing $\sigma$ by $s$:

\begin{definition}[Confidence and prediction intervals]
At a new $x_0$, with $\hat y_0 = b_0 + b_1x_0$ and $t^* = t^*_{n-2}$:
\begin{align*}
  \text{CI for the mean response:}\quad &
  \hat y_0 \pm t^*\, s\sqrt{\frac1n + \frac{(x_0-\xbar)^2}{S_{xx}}},\\
  \text{Prediction interval for a new $y$:}\quad &
  \hat y_0 \pm t^*\, s\sqrt{1 + \frac1n + \frac{(x_0-\xbar)^2}{S_{xx}}}.
\end{align*}
\end{definition}

For $x_0=4$: $(x_0-\xbar)^2/S_{xx} = 0.25/17.5 = 0.0143$, and $1/n = 0.1667$.
\begin{align*}
  \text{mean:}\quad & 6.6 \pm 2.776\times 0.837\sqrt{0.1810}
     = 6.6 \pm 0.99 = [5.61,\ 7.59],\\
  \text{new upload:}\quad & 6.6 \pm 2.776\times 0.837\sqrt{1.1810}
     = 6.6 \pm 2.52 = [4.08,\ 9.12].
\end{align*}

\begin{figure}[H]
\centering
\begin{tikzpicture}
\begin{axis}[
  width=11cm, height=7cm, axis lines=left,
  xmin=0, xmax=7.6, ymin=-1, ymax=16,
  xlabel={file size $x$ (MB)}, ylabel={upload time $y$ (s)},
  legend pos=north west, legend style={font=\small, cells={anchor=west}},
]
\addplot[name path=pu, draw=quotientorange, thick, dashed, domain=0.5:7.5, samples=60] {1.8+1.2*x+2.7764*0.83666*sqrt(1+1/6+(x-3.5)^2/17.5)};
\addlegendentry{95\% prediction band}
\addplot[name path=pl, draw=quotientorange, thick, dashed, domain=0.5:7.5, samples=60, forget plot] {1.8+1.2*x-2.7764*0.83666*sqrt(1+1/6+(x-3.5)^2/17.5)};
\addplot[name path=cu, draw=sourceblue, thick, domain=0.5:7.5, samples=60] {1.8+1.2*x+2.7764*0.83666*sqrt(1/6+(x-3.5)^2/17.5)};
\addlegendentry{95\% confidence band (mean)}
\addplot[name path=cl, draw=sourceblue, thick, domain=0.5:7.5, samples=60, forget plot] {1.8+1.2*x-2.7764*0.83666*sqrt(1/6+(x-3.5)^2/17.5)};
\addplot[quotientorange, fill opacity=0.12, forget plot] fill between[of=pu and pl];
\addplot[sourceblue, fill opacity=0.25, forget plot] fill between[of=cu and cl];
\addplot[black, very thick, domain=0.5:7.5, forget plot] {1.8+1.2*x};
\addplot[only marks, mark=*, mark size=2pt, black, forget plot] coordinates {(1,3) (2,5) (3,4) (4,7) (5,8) (6,9)};
\draw[dotted, gray] (axis cs:3.5,-1) -- (axis cs:3.5,6);
\node[gray, anchor=south west, font=\small] at (axis cs:3.5,-1) {$\xbar$};
\end{axis}
\end{tikzpicture}
\caption{Confidence band for the mean response (blue) and prediction band for a
single new upload (orange). Both are narrowest at $\xbar=3.5$ and widen with
distance from it. The prediction band is much wider because it includes the
noise of one upload.}
\label{fig:slr-bands}
\end{figure}

\begin{intuitionbox}
As $n\to\infty$ the confidence band shrinks to the true line: with enough data
you know the average exactly. The prediction band does not shrink below
$\pm t^* \sigma$: no amount of data predicts the jitter of one specific upload.
Dashboards that show a ``confidence band'' around a forecast are usually
showing the narrow one. If you want to know where tomorrow's single value will
land, you need the wide one.
\end{intuitionbox}

\begin{lstlisting}[language=Python]
n, x0 = len(x), 4.0
yhat = fit.intercept + fit.slope * x
s = np.sqrt(np.sum((y - yhat)**2) / (n - 2))      # 0.837
Sxx = np.sum((x - x.mean())**2)                   # 17.5
y0 = fit.intercept + fit.slope * x0               # 6.6
tc = stats.t.ppf(0.975, n - 2)                    # 2.776
se_mean = s * np.sqrt(1/n + (x0 - x.mean())**2 / Sxx)
se_pred = s * np.sqrt(1 + 1/n + (x0 - x.mean())**2 / Sxx)
print(y0 - tc*se_mean, y0 + tc*se_mean)           # 5.61 7.59
print(y0 - tc*se_pred, y0 + tc*se_pred)           # 4.08 9.12
\end{lstlisting}

\section{The assumptions: LINE}
\label{sec:slr-line}

The least squares line can always be computed. Its standard errors,
$p$-values and intervals rest on four assumptions, remembered as LINE.
Figure~\ref{fig:slr-line} lists each one, what breaks without it, and how to
check it.

\begin{figure}[H]
\centering
\begin{tikzpicture}[
  box/.style={draw, rounded corners, thick, minimum height=1.25cm,
    text width=2.8cm, align=center, font=\small},
  note/.style={text width=2.9cm, align=center, font=\footnotesize},
]
\node[box, draw=sourceblue, fill=sourceblue!8] (L) at (0,0) {\textbf{L}inearity\\$\E[y\given x]$ is a line};
\node[box, draw=targetgreen, fill=targetgreen!8] (I) at (3.2,0) {\textbf{I}ndependence\\each $\eps_i$ separate};
\node[box, draw=quotientorange, fill=quotientorange!8] (N) at (6.4,0) {\textbf{N}ormality\\$\eps\sim$ Normal};
\node[box, draw=structurepurple, fill=structurepurple!8] (E) at (9.6,0) {\textbf{E}qual variance\\same $\sigma$ for all $x$};
\node[note, below=0.25cm of L] (L2) {check: residual plot, no curve\\[2pt] fails $\Rightarrow$ slope is biased, predictions wrong};
\node[note, below=0.25cm of I] (I2) {check: study design; residuals vs.\ time\\[2pt] fails $\Rightarrow$ SEs too small};
\node[note, below=0.25cm of N] (N2) {check: histogram or QQ plot of residuals\\[2pt] fails $\Rightarrow$ matters for small $n$ and for PIs};
\node[note, below=0.25cm of E] (E2) {check: residual plot, no fan\\[2pt] fails $\Rightarrow$ SEs and intervals wrong};
\end{tikzpicture}
\caption{The four LINE assumptions, how to check each, and the cost of
violating it. Linearity is about the fit; the other three are about the error
bars.}
\label{fig:slr-line}
\end{figure}

The assumptions are not equally important.
\begin{itemize}
  \item \textbf{Linearity} matters most. If the mean is not a line, $b_1$
  estimates nothing meaningful.
  \item \textbf{Independence} is the one you cannot check from a scatter plot.
  Repeated measurements on the same user, or a time series where today
  resembles yesterday, make the effective sample size smaller than $n$ and the
  reported standard errors too small.
  \item \textbf{Normality} matters least for the slope. By the Central Limit
  Theorem $b_1 = \beta_1 + \sum w_i\eps_i$ is close to Normal for moderate $n$
  even when $\eps$ is not. It matters more for prediction intervals, which are
  about a single $\eps$.
  \item \textbf{Equal variance} affects every standard error. A fan-shaped
  residual plot means the intervals are too narrow in some places and too wide
  in others.
\end{itemize}

\section{Transformations and extrapolation}
\label{sec:slr-transform}

Many relationships are straight only on a log scale. Figure~\ref{fig:slr-loglog}
shows thirty animals' body mass and resting metabolic rate (simulated from
Kleiber's law). On the raw scale the points crowd into the corner and bend. On
log-log axes they fall on a line.

\begin{figure}[H]
\centering
\begin{tikzpicture}
\begin{axis}[
  name=left, width=6.3cm, height=5.6cm, axis lines=left,
  xmin=0, xmax=900, ymin=0, ymax=600,
  xlabel={body mass (kg)}, ylabel={metabolic rate (W)},
  title style={font=\small}, title={\small raw scale}, tick label style={font=\footnotesize},
  scaled ticks=false,
]
\addplot[sourceblue, thick, domain=0.1:880, samples=100] {2.9168*x^0.74695};
\addplot[only marks, mark=*, mark size=1.4pt] coordinates {(0.33,1.48) (102.01,98.42) (237.01,152.89) (189.28,195.67) (281.33,199.31) (32.88,32.23) (6.40,10.47) (318.33,228.03) (5.20,7.11) (1.79,6.40) (14.71,23.63) (8.34,15.69) (0.51,2.16) (600.62,231.09) (224.98,162.30) (579.30,530.58) (12.31,27.51) (9.59,13.55) (704.03,528.53) (2.16,5.90) (32.72,28.92) (5.82,8.74) (403.11,253.59) (90.71,65.52) (51.22,40.70) (31.95,28.76) (827.78,545.85) (1.12,2.29) (0.64,1.52) (1.52,7.53)};
\end{axis}
\begin{axis}[
  at={(left.east)}, anchor=west, xshift=0.6cm,
  width=6.3cm, height=5.6cm, axis lines=left,
  xmin=-0.7, xmax=3.1, ymin=0, ymax=3,
  xlabel={$\log_{10}$ mass}, ylabel={$\log_{10}$ rate},
  title style={font=\small}, title={\small log-log scale}, tick label style={font=\footnotesize},
]
\addplot[sourceblue, thick, domain=-0.6:3] {0.46490+0.74695*x};
\addplot[only marks, mark=*, mark size=1.4pt] coordinates {(-0.487,0.170) (2.009,1.993) (2.375,2.184) (2.277,2.292) (2.449,2.300) (1.517,1.508) (0.806,1.020) (2.503,2.358) (0.716,0.852) (0.252,0.806) (1.168,1.374) (0.921,1.196) (-0.291,0.334) (2.779,2.364) (2.352,2.210) (2.763,2.725) (1.090,1.440) (0.982,1.132) (2.848,2.723) (0.335,0.771) (1.515,1.461) (0.765,0.942) (2.605,2.404) (1.958,1.816) (1.709,1.610) (1.505,1.459) (2.918,2.737) (0.050,0.359) (-0.196,0.182) (0.181,0.877)};
\node[anchor=north west, font=\footnotesize, sourceblue] at (axis cs:-0.6,2.9) {slope $0.747$};
\end{axis}
\end{tikzpicture}
\caption{Left: a power law on raw axes is curved, and its noise fans out. Right:
on log-log axes it is a straight line with even scatter. The fitted log-log
slope $0.747$ is the exponent; the blue curve on the left is
$\hat y = 2.92\,x^{0.747}$.}
\label{fig:slr-loglog}
\end{figure}

If $y = a x^{b}$ times multiplicative noise, then
\[
  \log y = \log a + b\log x + \text{noise},
\]
which is the simple linear regression model in the variables $(\log x,\log y)$.
The fit gives $\log_{10} \hat y = 0.465 + 0.747\log_{10} x$, so
$\hat y = 10^{0.465}x^{0.747} = 2.92\,x^{0.747}$.

\textbf{Interpreting log-log slopes.} The slope is an \emph{elasticity}: a
$1\%$ increase in $x$ goes with about a $b\%$ increase in $y$. Doubling body mass
multiplies metabolic rate by $2^{0.747} \approx 1.68$. Other common forms:
\begin{itemize}
  \item $\log y$ on $x$ (log-linear): one unit of $x$ multiplies $y$ by
  $e^{b_1}$; for small $b_1$, that is about a $100b_1\%$ change. Natural for
  growth and for latencies, which are right-skewed.
  \item $y$ on $\log x$ (linear-log): multiplying $x$ by $10$ adds
  $b_1$ to $y$ (for $\log_{10}$). Natural for diminishing returns.
\end{itemize}
A log transform often fixes curvature and a fan shape at the same time, as
here. Use it when $y$ is positive and its spread grows with its level.

\begin{trapbox}
A line fitted to $\log y$ predicts the mean of $\log y$. Back-transforming
gives $10^{\widehat{\log y}}$, which estimates the \emph{median} of $y$ (if the
log-scale noise is symmetric), not the mean. The mean of $y$ is larger.
\end{trapbox}

\textbf{Extrapolation.} A regression line describes the data inside the range
of $x$ you observed. Outside that range you have no evidence the line
continues. Figure~\ref{fig:slr-extrap} shows eight measurements for
$x\in[0.5,4]$ that a line fits beautifully ($R^2 = 0.96$). The true curve
levels off. At $x=12$ the line predicts $21.3$; the truth is $9.9$.

\begin{figure}[H]
\centering
\begin{tikzpicture}
\begin{axis}[
  width=11cm, height=6.2cm, axis lines=left,
  xmin=0, xmax=13, ymin=0, ymax=24,
  xlabel={$x$}, ylabel={$y$}, clip=false,
  legend pos=north west, legend style={font=\small, cells={anchor=west}},
]
\fill[gray!12] (axis cs:0.5,0) rectangle (axis cs:4,24);
\addplot[sourceblue, very thick, domain=0:4] {1.3582+1.66054*x};
\addlegendentry{fitted line (data range)}
\addplot[sourceblue, very thick, dashed, domain=4:12] {1.3582+1.66054*x};
\addlegendentry{line extrapolated}
\addplot[targetgreen, very thick, domain=0:12, samples=80] {10*(1-exp(-0.35*x))};
\addlegendentry{true mean curve}
\addplot[only marks, mark=*, mark size=2pt, black, forget plot] coordinates {(0.50,1.81) (1.00,2.65) (1.50,4.18) (2.00,5.28) (2.50,5.63) (3.00,6.65) (3.50,6.96) (4.00,7.58)};
\node[font=\small, anchor=south] at (axis cs:2.25,0.4) {observed};
\addplot[only marks, mark=*, mark size=2pt, trapred, forget plot] coordinates {(12,21.28) (12,9.85)};
\node[trapred, anchor=west, font=\small] at (axis cs:12.15,21.28) {21.3};
\node[trapred, anchor=west, font=\small] at (axis cs:12.15,9.85) {9.9};
\end{axis}
\end{tikzpicture}
\caption{Extrapolation. Inside the shaded data range the line is an excellent
approximation. Beyond it the line and the truth diverge, and nothing in the
data could have warned you.}
\label{fig:slr-extrap}
\end{figure}

The prediction interval formula grows with $(x_0-\xbar)^2$, so it does widen
outside the data. But it only accounts for uncertainty in the slope and
intercept \emph{assuming the line is right}. It cannot account for the line
being wrong, which is the bigger risk far from the data. Capacity planning
(``latency grows linearly with load, so at $3\times$ load\ldots'') is the classic
victim: queues are linear at low utilisation and explode near saturation.

\section{Least squares is maximum likelihood}
\label{sec:slr-mle}

Section~\ref{sec:slr-ls} chose least squares because it is convenient. There
is a better reason. Assume the Normal-error model, $y_i \sim \Normal(\beta_0+\beta_1x_i,\
\sigma^2)$ independently. The probability density of the observed data, viewed
as a function of the parameters, is the \textbf{likelihood}:
\[
  L(\beta_0,\beta_1,\sigma)
  = \prod_{i=1}^n \frac{1}{\sqrt{2\pi}\,\sigma}
    \exp\!\left(-\frac{(y_i-\beta_0-\beta_1x_i)^2}{2\sigma^2}\right).
\]
Take logs:
\[
  \log L = -\frac{n}{2}\log(2\pi\sigma^2)
           - \frac{1}{2\sigma^2}\sum_{i=1}^n (y_i-\beta_0-\beta_1x_i)^2 .
\]
For any fixed $\sigma$, the only term involving $\beta_0,\beta_1$ is
$-\mathrm{SSE}/(2\sigma^2)$. Maximising $\log L$ is therefore the same as
minimising $\mathrm{SSE}$. The maximum likelihood estimates of slope and
intercept are exactly the least squares estimates.

Maximising over $\sigma$ as well gives $\hat\sigma^2_{\text{MLE}} =
\mathrm{SSE}/n$. That is biased low; the unbiased version divides by $n-2$, the
$s^2$ of Section~\ref{sec:slr-inference}. Same story as $n$ versus $n-1$ for
the sample variance.

\begin{connectionbox}
This equivalence is why ``mean squared error loss'' is the default in machine
learning regression. Training a neural network with MSE loss is maximum
likelihood under the assumption that targets are the network's output plus
Gaussian noise. Swap the noise model and the loss changes: Laplace noise gives
absolute error (L1) loss; Bernoulli outcomes give cross-entropy, which is
logistic regression (Chapter~13). Chapter~14 develops maximum likelihood in
general, and Chapter~15 returns to the loss-function view.
\end{connectionbox}

Upload time surely depends on more than file size: network type, server
region, time of day. Chapter~13 puts several predictors into one regression
and asks what each slope means when the others are held fixed.

\section{Summary}
\label{sec:slr-summary}

\begin{center}
\renewcommand{\arraystretch}{1.25}
\begin{tabular}{p{4.1cm}p{8.4cm}}
\toprule
Model & $y_i = \beta_0+\beta_1x_i+\eps_i$, errors independent, mean $0$, variance $\sigma^2$ \\
Residual & $e_i = y_i-\hat y_i$, the vertical distance to the line \\
Least squares slope & $b_1 = S_{xy}/S_{xx} = r\,s_y/s_x$ \\
Least squares intercept & $b_0 = \ybar - b_1\xbar$; line passes through $(\xbar,\ybar)$ \\
Normal equations & $\sum e_i = 0$, $\sum x_ie_i = 0$ \\
Decomposition & $\mathrm{SST} = \mathrm{SSR}+\mathrm{SSE}$; $R^2 = \mathrm{SSR}/\mathrm{SST} = r^2$ \\
Regression to the mean & $\hat z_y = r\,z_x$; $y$-on-$x$ and $x$-on-$y$ lines differ unless $\abs{r}=1$ \\
Residual SE & $s = \sqrt{\mathrm{SSE}/(n-2)}$ \\
Slope inference & $\SE(b_1)=s/\sqrt{S_{xx}}$; $(b_1-\beta_1)/\SE(b_1)\sim t_{n-2}$ \\
Mean response CI & $\hat y_0\pm t^*s\sqrt{1/n+(x_0-\xbar)^2/S_{xx}}$ \\
Prediction interval & $\hat y_0\pm t^*s\sqrt{1+1/n+(x_0-\xbar)^2/S_{xx}}$ \\
Diagnostics & residual plot: curve $\Rightarrow$ nonlinearity; fan $\Rightarrow$ unequal variance; high leverage $+$ off trend $\Rightarrow$ influential \\
Assumptions & LINE: Linearity, Independence, Normality, Equal variance \\
Log-log slope & elasticity: $1\%$ more $x$ $\approx$ $b_1\%$ more $y$ \\
Likelihood view & Gaussian noise $\Rightarrow$ least squares $=$ maximum likelihood \\
\bottomrule
\end{tabular}
\end{center}

\section{Exercises}
\label{sec:slr-exercises}

\begin{exercise}
Five points: $x = 0,1,2,3,4$ and $y = 1,3,2,5,4$. Compute $\xbar$, $\ybar$,
$S_{xx}$, $S_{xy}$, $S_{yy}$, the least squares line, $r$ and $R^2$.
\end{exercise}

\begin{exercise}
For the data in Exercise 1, compute $\mathrm{SSE}$, $s$, $\SE(b_1)$ and the
$t$ statistic for $H_0\colon\beta_1=0$. With $t^*_3 = 3.182$, is the slope
significant at the $5\%$ level? What is the 95\% CI for $\beta_1$?
\end{exercise}

\begin{exercise}
A study reports $\xbar = 50$, $s_x = 10$, $\ybar = 20$, $s_y = 5$, $r = 0.6$.
Find the regression line of $y$ on $x$ and the regression line of $x$ on $y$.
Predict $y$ when $x = 70$.
\end{exercise}

\begin{exercise}
Midterm and final exam scores have correlation $r=0.5$. A student scores
$2$ SDs above the mean on the midterm. Predict how many SDs above the mean she
will score on the final. A student who scored $2$ SDs above the mean on the
final: predict her midterm in SD units. Why is neither answer $2$?
\end{exercise}

\begin{exercise}
Prove from the normal equations that the residuals of a least squares fit with
an intercept satisfy $\sum_i \hat y_i e_i = 0$.
\end{exercise}

\begin{exercise}
In the upload example, $x$ is converted from MB to KB ($1$ MB $=1000$ KB) and
$y$ from seconds to milliseconds. What happens to $b_1$, $b_0$, $R^2$ and the
$t$ statistic for the slope?
\end{exercise}

\begin{exercise}
For the upload data, compute the 95\% confidence interval for the mean response
and the 95\% prediction interval at $x_0 = 7$ MB. Why are both wider than at
$x_0 = 4$? What extra caveat applies at $x_0=7$?
\end{exercise}

\begin{exercise}
A log-log regression of house price on floor area gives slope $0.8$. By what
factor does the predicted price change when floor area doubles? When it grows
by $10\%$?
\end{exercise}

\begin{exercise}
You want to estimate how request latency depends on payload size, and you can
choose $20$ payload sizes in the range $1$ to $100$ KB. Using
$\SD(b_1) = \sigma/\sqrt{S_{xx}}$, which choice of sizes gives the most precise
slope? Why might you not use that choice in practice?
\end{exercise}

\begin{exercise}
Show that under Gaussian errors, maximising the log-likelihood over $\sigma^2$
(with $\beta_0,\beta_1$ at their least squares values) gives
$\hat\sigma^2 = \mathrm{SSE}/n$.
\end{exercise}

\subsection*{Answers}

\begin{enumerate}[leftmargin=*]
  \item $\xbar=2$, $\ybar=3$, $S_{xx}=10$, $S_{xy}=8$, $S_{yy}=10$.
  $b_1=0.8$, $b_0=3-0.8\times2=1.4$, so $\hat y = 1.4+0.8x$.
  $r = 8/\sqrt{100}=0.8$, $R^2=0.64$.
  \item $\mathrm{SSE} = S_{yy}-b_1S_{xy} = 10-6.4=3.6$; $s=\sqrt{3.6/3}=1.095$;
  $\SE(b_1)=1.095/\sqrt{10}=0.346$; $t=0.8/0.346=2.31<3.182$. Not significant
  ($p=0.10$). CI: $0.8\pm3.182\times0.346 = [-0.30,\ 1.90]$, which contains $0$.
  \item $y$ on $x$: $b_1 = 0.6\times5/10=0.3$, $b_0=20-0.3\times50=5$, so
  $\hat y=5+0.3x$. $x$ on $y$: $c_1=0.6\times10/5=1.2$, $c_0=50-1.2\times20=26$,
  so $\hat x=26+1.2y$. At $x=70$: $\hat y = 26$.
  \item Final: $0.5\times2 = 1$ SD above. Midterm given final: also $1$ SD
  above. Each prediction regresses toward the mean because $\abs{r}<1$; the
  extreme score is partly luck, which does not carry over.
  \item $\hat y_i = b_0+b_1x_i$, so
  $\sum\hat y_ie_i = b_0\sum e_i + b_1\sum x_ie_i = 0+0$.
  \item $b_1$ becomes $1.2\times1000/1000=1.2$ ms/KB (unchanged numerically
  here, since both scales changed by $1000$). $b_0$ becomes $1800$ ms. $R^2$
  and $t$ are unchanged: they are unit-free.
  \item $\hat y_0=10.2$. $(7-3.5)^2/17.5=0.7$. Mean CI:
  $10.2\pm2.776\times0.837\sqrt{0.8667} = [8.04,\ 12.36]$. PI:
  $10.2\pm2.776\times0.837\sqrt{1.8667}=[7.03,\ 13.37]$. Wider because $x_0$ is
  farther from $\xbar$. And $x_0=7$ is outside the observed range $[1,6]$, so it
  is an extrapolation: the intervals assume the line still holds there.
  \item Doubling: $2^{0.8}\approx1.74$, a $74\%$ increase. $10\%$ more area:
  $1.1^{0.8}\approx1.079$, about $7.9\%$ more.
  \item Put $10$ sizes at $1$ KB and $10$ at $100$ KB; this maximises
  $S_{xx}$. But with only two distinct $x$ values you cannot see curvature, so
  you cannot check linearity. Spreading some points across the range is safer.
  \item $\partial\log L/\partial\sigma^2 = -n/(2\sigma^2) +
  \mathrm{SSE}/(2\sigma^4)=0$ gives $\sigma^2 = \mathrm{SSE}/n$.
\end{enumerate}
